\chapter{Introduction}
\label{ch:introduction}

\section{Why automated reaction-path search}

Reaction networks are the map of a chemical system: which elementary steps
connect which intermediates, and how high the barriers along the way are.
Mapping them by hand is slow and biased -- any path an expert does not think
of is a path that is never computed.  Automated exploration methods, such as
the artificial force induced reaction (AFIR) method \cite{maeda2016afir},
remove the bias: an artificial force term is added to the potential energy
surface (PES) to push fragments together or apart, and the modified surface
is minimised to obtain approximate reaction paths.

The cost is the bottleneck.  Every trial path needs hundreds of energy and
gradient evaluations, traditionally from density functional theory (DFT), and
an exploration involves thousands of trials.  Two levers are available to
make this affordable: better machinery (parallel scheduling, job daemons)
and cheaper potential energy evaluations.  \jns{} attacks the second lever.

\section{Machine-learning potentials, and what they cannot promise}

Machine-learning potentials (MLPs) fit the PES of a family of molecules with
neural networks, reproducing near-DFT energies and forces at a small
fraction of the cost.  Modern architectures -- MACE's equivariant message
passing \cite{batatia2022mace} and the MACE-MP-0 materials foundation model
\cite{batatia2025mpa}, or AIMNet2 for molecular and elemental-organic
chemistry \cite{anstine2025aimnet2} -- are accurate on chemistry they were
trained for.

The trap is that reaction-path search deliberately leaves the training
distribution.  AFIR pushes geometries into compressed and stretched
structures -- exactly the high-energy, bond-forming and bond-breaking
configurations that near-equilibrium training sets under-represent.  A
model that is accurate to 1 kcal/mol on equilibrium molecules can be wrong
by an order of magnitude on a strained saddle region, and the error is
systematic and silent: nothing in the output says ``this number is outside
my domain''.  The Passerini controversy studies
\cite{staub2023,staub2025} make the same point from the other direction: a
carefully iterated potential was needed before a mechanism question could be
settled.

\section{The design idea: uncertainty as a run-time decision}

\jns{} takes the position that a fast prediction is only useful together
with a statement of its reliability, and that this statement should be
\emph{measured}, not assumed.  Three ingredients implement this:

\begin{enumerate}
  \item \textbf{A committee of different architecture families.}  Two MACE
    models and one AIMNet2 model evaluate every geometry; their spread is a
    signal of how contested the prediction is.  Members from a single family
    would agree with each other almost perfectly and carry no signal.
  \item \textbf{A calibrated threshold.}  The disagreement score is
    calibrated against reference errors, so that the probability of a
    dangerous geometry slipping below the threshold is bounded by
    construction (Chapter \ref{ch:uncertainty}).
  \item \textbf{A hard fallback.}  When the threshold is exceeded, the
    evaluation goes to ORCA at the production level instead of the MLP.  The
    fallback is real quantum chemistry, not an extrapolated correction.
\end{enumerate}

The framework itself does not run your exploration for you: it supplies the
calculators and the integration glue, you run the engine.  The exploration
engine is SCINE Chemoton \cite{unsleber2022chemoton,weymuth2024scine} with
its job daemon Puffin \cite{bensberg2025puffin}, both open source; the whole
stack -- SCINE, MACE, AIMNet2, ASE \cite{larsen2017ase} -- is publicly
available, with ORCA \cite{neese2022orca} as the academic reference program.

\section{Relation to existing work}

Driving reaction-path exploration with MLPs is an active line of work.  The
GRRM-based studies of Staub et al.\ \cite{staub2023,staub2025} train
SpookyNet-family potentials for targeted systems; the SCINE group's Parrot
module \cite{eckhoff2025parrot} integrates lifelong-learning potentials into
the same engine family used here.  \jns{} occupies the adjacent position of
a lightweight, fully open adaptation route: it reuses public foundation
models instead of training a general-purpose potential from scratch, adapts
them to a reaction family with a few hundred labelled frames, and adds the
run-time uncertainty layer described above.  Its purpose is not to replace
the existing routes but to lower the entry cost of running one.

\section{How this manual is organised}

Chapters \ref{ch:installation}--\ref{ch:getting-started} get you running
quickly.  Chapters \ref{ch:architecture}--\ref{ch:uncertainty} explain the
design: the three-layer architecture, the calculators, the SCINE
integration, and the uncertainty machinery.  Chapter
\ref{ch:active-learning} documents the active-learning workflow and
Chapter \ref{ch:scripts} the utility scripts.  Chapters
\ref{ch:configuration} and \ref{ch:development} are reference material for
configuration fields and for development practices.  The appendices collect
constants, environment variables and terminology.
